% GENERATED FILE. Edit canonical lessons, then run npm run generate:books.
% canonical-source-hash: fnv1a64:beae2b71f7590126
% canonical-lessons: ZH-C176-qiaomen, ZH-C176-linju, ZH-C176-chaojia, ZH-C176-dao-topple, ZH-C176-zhaohuo

\chapter{A Knock at the Door}
\label{ch:zh-a-knock-at-the-door}
\hlchaptermodality{176}

\begin{chapteropening}
\textbf{By the end of this chapter:} \emph{I can tell what happened at home: a knock at the door, the neighbours quarrelling, something toppling over, a fire.}

Complete the last lesson of chapter 176: I can tell what happened at home: a knock at the door, the neighbours quarrelling, something toppling over, a fire.
\end{chapteropening}

\section[qiāomén]{qiāomén — to knock at the door}

\label{lesson:ZH-C176-qiaomen}



\begin{quote}
Before the new one: say the Mandarin for a cupboard, then the Mandarin for a bookshelf.
\end{quote}

\subsection*{You'll want to know: qiāomén}
\emph{qiāomén} — \textquotedblleft{}to knock at the door\textquotedblright{}.

\begin{quote}
\emph{zuótiān wǎnshang yǒu rén qiāomén} — Somebody knocked at the door last night.
\end{quote}

Its characters come later; for now it is a word for the ear.

\begin{cousinweb}
\emph{qiāomén} is built from \emph{qiāo} 'to strike' and \zh{门} 'door', both inherited from Middle Chinese, \emph{khaew} and \emph{mwon} (Guangyun, in Baxter's transcription).
\end{cousinweb}

\begin{grammarlens}[title={What you've built}]
One more word for talking about the past.
\end{grammarlens}

\subsection*{Your turn}
\begin{itemize}
\raggedright
  \item \emph{Say it:} \emph{qiāomén}
  \item \emph{Say it:} \emph{qiāomén}, once more
  \item \emph{Say it:} \emph{zuótiān wǎnshang yǒu rén qiāomén} — Somebody knocked at the door last night.
  \item \emph{Recall:} say the Mandarin for a cupboard, then the Mandarin for a bookshelf, then say \emph{qiāomén} again
\end{itemize}

\begin{tcolorbox}[breakable,colback=teal!4,colframe=teal!35!black,title={Before you move on}]
What does qiāomén mean? (\textquotedblleft{}To knock at the door\textquotedblright{}.) Say it once more.
\end{tcolorbox}
\section[línjū]{línjū — neighbour}

\label{lesson:ZH-C176-linju}



\begin{quote}
Before the new one: say the Mandarin for a bookshelf, then the Mandarin for to knock at the door.
\end{quote}

\subsection*{You'll want to know: línjū}
\emph{línjū} — \textquotedblleft{}neighbour\textquotedblright{}.

\begin{quote}
\emph{qiāomén de rén shì wǒ de línjū} — The one knocking was my neighbour.
\end{quote}

Its characters come later; for now it is a word for the ear.

\begin{cousinweb}
\emph{línjū} is built from \emph{lín} 'near' and \emph{jū} 'to dwell', both inherited from Middle Chinese, \emph{lin} and \emph{kjo} (Guangyun, in Baxter's transcription): the one who dwells near. After an older book, the dictionary also counts five households together as one \emph{lín}.
\end{cousinweb}

\begin{grammarlens}[title={What you've built}]
One more word for talking about the past.
\end{grammarlens}

\subsection*{Your turn}
\begin{itemize}
\raggedright
  \item \emph{Say it:} \emph{línjū}
  \item \emph{Say it:} \emph{línjū}, once more
  \item \emph{Say it:} \emph{qiāomén de rén shì wǒ de línjū} — The one knocking was my neighbour.
  \item \emph{Recall:} say the Mandarin for a bookshelf, then the Mandarin for to knock at the door, then say \emph{línjū} again
\end{itemize}

\begin{tcolorbox}[breakable,colback=teal!4,colframe=teal!35!black,title={Before you move on}]
What does línjū mean? (\textquotedblleft{}Neighbour\textquotedblright{}.) Say it once more.
\end{tcolorbox}
\section[chǎojià]{chǎojià — to quarrel; to have a row}

\label{lesson:ZH-C176-chaojia}



\begin{quote}
Before the new one: say the Mandarin for to knock at the door, then the Mandarin for a neighbour.
\end{quote}

\subsection*{You'll want to know: chǎojià}
\emph{chǎojià} — \textquotedblleft{}to quarrel; to have a row\textquotedblright{}.

\begin{quote}
\emph{tā chángcháng gēn línjū chǎojià} — He's forever quarrelling with the neighbours.
\end{quote}

Its characters come later; for now it is a word for the ear.

\begin{cousinweb}
\emph{chǎojià} is built from \zh{吵} 'noisy' and \emph{jià}. The first is inherited from Middle Chinese \emph{tsrhaewX} (Guangyun, in Baxter's transcription), glossed 'a sound'; the second is the word for a frame, and how it came to mean a fight is not settled.
\end{cousinweb}

\begin{grammarlens}[title={What you've built}]
One more word for talking about the past.
\end{grammarlens}

\subsection*{Your turn}
\begin{itemize}
\raggedright
  \item \emph{Say it:} \emph{chǎojià}
  \item \emph{Say it:} \emph{chǎojià}, once more
  \item \emph{Say it:} \emph{tā chángcháng gēn línjū chǎojià} — He's forever quarrelling with the neighbours.
  \item \emph{Recall:} say the Mandarin for to knock at the door, then the Mandarin for a neighbour, then say \emph{chǎojià} again
\end{itemize}

\begin{tcolorbox}[breakable,colback=teal!4,colframe=teal!35!black,title={Before you move on}]
What does chǎojià mean? (\textquotedblleft{}To quarrel; to have a row\textquotedblright{}.) Say it once more.
\end{tcolorbox}
\section[dǎo]{dǎo — to topple over; to fall down (of a thing)}

\label{lesson:ZH-C176-dao-topple}



\begin{quote}
Before the new one: say the Mandarin for a neighbour, then the Mandarin for to quarrel.
\end{quote}

\subsection*{You'll want to know: dǎo}
\emph{dǎo} — \textquotedblleft{}to topple over; to fall down (of a thing)\textquotedblright{}.

\begin{quote}
\emph{tā de zìxíngchē dǎo le} — His bike fell over.
\end{quote}

Its characters come later; for now it is a word for the ear.

\begin{cousinweb}
\emph{dǎo} is inherited from Middle Chinese \emph{tawX} (Guangyun, in Baxter's transcription), glossed 'to fall forward'.
\end{cousinweb}

\begin{grammarlens}[title={What you've built}]
One more word for talking about the past.
\end{grammarlens}

\subsection*{Your turn}
\begin{itemize}
\raggedright
  \item \emph{Say it:} \emph{dǎo}
  \item \emph{Say it:} \emph{dǎo}, once more
  \item \emph{Say it:} \emph{tā de zìxíngchē dǎo le} — His bike fell over.
  \item \emph{Recall:} say the Mandarin for a neighbour, then the Mandarin for to quarrel, then say \emph{dǎo} again
\end{itemize}

\begin{tcolorbox}[breakable,colback=teal!4,colframe=teal!35!black,title={Before you move on}]
What does dǎo mean? (\textquotedblleft{}To topple over; to fall down (of a thing)\textquotedblright{}.) Say it once more.
\end{tcolorbox}
\section[zháohuǒ]{zháohuǒ — to catch fire}

\label{lesson:ZH-C176-zhaohuo}



\begin{quote}
Before the new one: say the Mandarin for to quarrel, then the Mandarin for to topple over.
\end{quote}

\subsection*{You'll want to know: zháohuǒ}
\emph{zháohuǒ} — \textquotedblleft{}to catch fire\textquotedblright{}.

\begin{quote}
\emph{qiántiān línjū jiā zháohuǒ le} — The neighbours' house caught fire the day before yesterday.
\end{quote}

Its characters come later; for now it is a word for the ear.

\begin{cousinweb}
\emph{zháohuǒ} is built from \emph{zháo} 'to attach to' and \zh{火} 'fire', both inherited from Middle Chinese, \emph{drjak} and \emph{xwaX} (Guangyun, in Baxter's transcription): fire takes hold.
\end{cousinweb}

\begin{grammarlens}[title={What you've built}]
One more word for talking about the past.
\end{grammarlens}

\subsection*{Your turn}
\begin{itemize}
\raggedright
  \item \emph{Say it:} \emph{zháohuǒ}
  \item \emph{Say it:} \emph{zháohuǒ}, once more
  \item \emph{Say it:} \emph{qiántiān línjū jiā zháohuǒ le} — The neighbours' house caught fire the day before yesterday.
  \item \emph{Recall:} say the Mandarin for to quarrel, then the Mandarin for to topple over, then say \emph{zháohuǒ} again
\end{itemize}

\begin{tcolorbox}[breakable,colback=teal!4,colframe=teal!35!black,title={Before you move on}]
What does zháohuǒ mean? (\textquotedblleft{}To catch fire\textquotedblright{}.) Say it once more.
\end{tcolorbox}
